Fix \(n\), and first let \(P\in\mathcal P_n(\Lambda_n)\). By \(\mathrm{def:triangular\text{-}profile\text{-}union}\), there is a declared triple with exponent \(\kappa\) for which \(P\in\mathcal P_{\kappa}^{\mathrm{base}}\). Hence the Lipschitz, bounded-trace, Gaussian-error, polynomial-lower-mass, and independent-sampling conditions hold by \(\mathrm{def:base\text{-}thickness\text{-}class}\) and the corresponding assumptions. In particular, \(\mathrm{prop:observed\text{-}law\text{-}identification}\) gives, for every \(x\in\mathcal B\) and \(t\in\{0,1\}\),
\[
x\in\mathcal X_{t,P}
\qquad\text{and}\qquad
\theta_{t,P}(x)=g_{t,P}(x). \tag{1}
\]
Thus the object covered below is the identified trace contrast \(\tau_P=\theta_{1,P}-\theta_{0,P}\), while \(\widehat\tau_n\) remains the explicit sample map in \(\mathrm{def:mass\text{-}adaptive\text{-}estimator}\).

For fixed scores \(X_1,\ldots,X_n\), write \(\mathscr S_n\) for the nonempty disk-pattern family in \(\mathrm{lem:disk\text{-}pattern\text{-}gaussian\text{-}control}\). That lemma and the definition of \(M_n\) give
\[
|\mathscr S_n|\le M_n. \tag{2}
\]
Set
\[
z_{n,\gamma}=\sqrt{2\log(2M_n/\gamma)}.
\]
The conditional Gaussian inequality in the same lemma, followed by (2), yields
\[
\begin{aligned}
&P\left\{
\max_{S\in\mathscr S_n}
\frac{\left|\sum_{i\in S}\{Y_i-g_{T_i,P}(X_i)\}\right|}
{\sigma\sqrt{|S|}}>z_{n,\gamma}
\ \middle|\ X_1,\ldots,X_n\right\}\\
&\hspace{35mm}\le 2|\mathscr S_n|e^{-z_{n,\gamma}^2/2}
\le\gamma.
\end{aligned}\tag{3}
\]
The selected bandwidth \(\widehat h_{t,n}(x)\) is a function only of the scores and the known geometry. Consequently, conditional on the scores, every nonempty index set selected by \((x,t)\) belongs to \(\mathscr S_n\), and (3) holds simultaneously for all such selections.

Fix \(x\in\mathcal B\) and \(t\in\{0,1\}\), abbreviate \(\widehat h=\widehat h_{t,n}(x)\) and \(N=N_{t,n}(x,\widehat h)\), and suppose first that \(N>0\). Apart from a \(P\)-null score event, every selected \(X_i\) lies in \(\mathcal X_{t,P}\). Since \(d(X_i,x)\le\widehat h\le h_0\), assumption \(\mathrm{ass:lipschitz\text{-}side\text{-}means}\) and (1) imply
\[
\left|\frac1N\sum_{i:X_i\in\mathcal A_t,\ d(X_i,x)\le\widehat h}
g_{t,P}(X_i)-\theta_{t,P}(x)\right|
\le L\widehat h. \tag{4}
\]
On the complementary event to that in (3), the stochastic part satisfies
\[
\left|\frac1N\sum_{i:X_i\in\mathcal A_t,\ d(X_i,x)\le\widehat h}
\{Y_i-g_{t,P}(X_i)\}\right|
\le \sigma\sqrt{\frac{2\log(2M_n/\gamma)}{N}}. \tag{5}
\]
Combining \(\mathrm{def:local\text{-}constant\text{-}side\text{-}fit}\), (4), and (5) gives
\[
|\widehat g_{t,n}(x;\widehat h)-\theta_{t,P}(x)|\le B_{t,n}(x). \tag{6}
\]
If \(N=0\), the empty sum in \(\mathrm{def:local\text{-}constant\text{-}side\text{-}fit}\) is zero. Assumption \(\mathrm{ass:bounded\text{-}side\text{-}traces}\) then gives
\[
|\widehat g_{t,n}(x;\widehat h)-\theta_{t,P}(x)|
=|\theta_{t,P}(x)|\le L=B_{t,n}(x), \tag{7}
\]
so (6) remains true. By the definitions of \(\tau_P\) and \(\widehat\tau_n\), the triangle inequality applied to (6) for \(t=0,1\) yields, simultaneously in \(x\),
\[
|\widehat\tau_n(x)-\tau_P(x)|
\le B_{0,n}(x)+B_{1,n}(x). \tag{8}
\]
Equations (3) and (8) prove conditional coverage at least \(1-\gamma\), hence unconditional coverage at least \(1-\gamma\). The argument is uniform in \(P\in\mathcal P_n(\Lambda_n)\), proving the displayed infimum guarantee.

The band is data measurable: the grid is finite; every count and empirical mean is measurable; \(\widehat{\mathcal H}_{t,n}(x)\) is therefore measurable; and \(\mathrm{def:count\text{-}adaptive\text{-}bandwidth}\) specifies both a deterministic numerical tie break and a deterministic empty-set fallback. Its two endpoints are consequently measurable functions of the data and the known geometry.

We next prove the expected-width assertion directly. Fix \(\kappa\in(2,\overline\kappa]\), let \(P\in\mathcal P_{\kappa}^{\mathrm{perv}}\), and put
\[
s_{n,\kappa}=\left(\frac{\log n}{n}\right)^{1/(\kappa+2)}
=r_{n,\kappa}^{\mathrm{perv}}. \tag{9}
\]
By \(\mathrm{def:pervasive\text{-}thinning\text{-}class}\), \(P\in\mathcal P_{\kappa}^{\mathrm{base}}\), so the lower-mass condition applies. For all sufficiently large \(n\), uniformly over \(\kappa\in(2,\overline\kappa]\), the dyadic grid contains an \(h_\star=h_{\star,n,\kappa}\) satisfying
\[
s_{n,\kappa}\le h_\star<2s_{n,\kappa}<h_0. \tag{10}
\]
Indeed, \(s_{n,\kappa}\to0\) uniformly at its slowest endpoint \(\overline\kappa\), while \(ns_{n,\kappa}\to\infty\) uniformly at its fastest endpoint, so \(s_{n,\kappa}\) lies between the smallest and largest grid elements; adjacent grid elements differ by a factor of two.

Let \(\mathcal G_\star\) be a maximal \(h_\star/4\)-separated subset of the compact boundary. It is an \(h_\star/4\)-net, and assumption \(\mathrm{ass:global\text{-}boundary\text{-}entropy}\) gives
\[
|\mathcal G_\star|\le 4C_{\mathcal B}h_\star^{-1}. \tag{11}
\]
For \(y\in\mathcal G_\star\) and \(t\in\{0,1\}\), define the auxiliary count
\[
Z_{t,y}=\sum_{i=1}^n
\mathbf 1\{X_i\in\mathcal A_t,\ d(X_i,y)\le h_\star/2\}.
\]
Assumption \(\mathrm{ass:iid\text{-}sampling}\) makes \(Z_{t,y}\) binomial with mean
\[
\mu_{t,y}=nq_{t,P}(y,h_\star/2)
\ge nc_m(h_\star/2)^\kappa, \tag{12}
\]
where the inequality is \(\mathrm{ass:polynomial\text{-}lower\text{-}mass}\). For completeness, if \(Z\) is binomial with mean \(\mu=nq\), Markov's inequality at \(\lambda=\log2\), together with \(1-u\le e^{-u}\), gives
\[
\begin{aligned}
P(Z\le\mu/2)
&\le e^{\lambda\mu/2}\mathbb E e^{-\lambda Z}\\
&=e^{\lambda\mu/2}(1-q+qe^{-\lambda})^n\\
&\le\exp\{\mu(\lambda/2-1+e^{-\lambda})\}
\le e^{-\mu/8}.
\end{aligned}\tag{13}
\]
The last inequality follows from \((\log2)/2-1/2<-1/8\). Define
\[
\mathcal E_{n,\kappa}
=\bigcap_{t\in\{0,1\}}\bigcap_{y\in\mathcal G_\star}
\{Z_{t,y}\ge\mu_{t,y}/2\}.
\]
Equations (11)--(13) and a union bound show, with constants uniform in \(P\) and \(\kappa\),
\[
P(\mathcal E_{n,\kappa}^c)
\le C h_\star^{-1}\exp\{-cnh_\star^\kappa\}. \tag{14}
\]

For every \(x\in\mathcal B\), choose \(y\in\mathcal G_\star\) with \(d(x,y)<h_\star/4\). The triangle inequality gives
\[
\{z:d(z,y)\le h_\star/2\}
\subseteq\{z:d(z,x)\le h_\star\}.
\]
Thus, on \(\mathcal E_{n,\kappa}\), (12) implies simultaneously in \(x\) and \(t\) that
\[
N_{t,n}(x,h_\star)
\ge \frac{c_m}{2^{\kappa+1}}nh_\star^\kappa
\ge c_0nh_\star^\kappa, \qquad
c_0:=c_m2^{-(\overline\kappa+1)}>0. \tag{15}
\]
Moreover, by (9)--(10),
\[
\frac{nh_\star^\kappa}{\log(e+n/h_\star)}
\ge c\left(\frac{n}{\log n}\right)^{2/(\kappa+2)}
\ge c\left(\frac{n}{\log n}\right)^{2/(\overline\kappa+2)}
\longrightarrow\infty. \tag{16}
\]
It follows from (15)--(16) that, for every sufficiently large \(n\),
\[
N_{t,n}(x,h_\star)\ge8\log(e+nh_\star^{-1})
\]
simultaneously in \(x,t,P,\kappa\). Hence \(h_\star\in\widehat{\mathcal H}_{t,n}(x)\) on \(\mathcal E_{n,\kappa}\).

For an admissible \(h\), set
\[
Q_{t,n}(x,h)=Lh+\sigma
\sqrt{\frac{\log(e+nh^{-1})}{N_{t,n}(x,h)}}.
\]
The minimizing property in \(\mathrm{def:count\text{-}adaptive\text{-}bandwidth}\), (10), and (15) give on \(\mathcal E_{n,\kappa}\)
\[
\begin{aligned}
Q_{t,n}(x,\widehat h_{t,n}(x))
&\le Q_{t,n}(x,h_\star)\\
&\le Lh_\star+C\sigma
\sqrt{\frac{\log n}{nh_\star^\kappa}}\\
&\le C s_{n,\kappa},
\end{aligned}\tag{17}
\]
because the two terms in the middle line balance exactly at (9):
\[
\sqrt{\frac{\log n}{ns_{n,\kappa}^\kappa}}
=s_{n,\kappa}. \tag{18}
\]
Also \(|\mathcal H_n|\le2+\log_2 n\) and \(\sum_{j=0}^3\binom nj\le Cn^3\), so, for fixed \(\gamma\in(0,1)\),
\[
\log(2M_n/\gamma)\le C_\gamma\log n. \tag{19}
\]
For every \(h\in\mathcal H_n\), the condition \(h\le h_0<1\) gives \(\log(e+nh^{-1})\ge\log n\). Combining this fact with (17)--(19) yields
\[
\sup_{x\in\mathcal B}\max_{t\in\{0,1\}}B_{t,n}(x)
\le Cs_{n,\kappa}
\qquad\text{on }\mathcal E_{n,\kappa}. \tag{20}
\]

Off \(\mathcal E_{n,\kappa}\), the definition of \(B_{t,n}\), \(N\ge1\) in its first branch, and \(h_0<1\) give the deterministic envelope
\[
\sup_{x\in\mathcal B}\max_t B_{t,n}(x)
\le L+\sigma\sqrt{2\log(2M_n/\gamma)}
\le C\sqrt{\log n}. \tag{21}
\]
Since the constructed interval is symmetric,
\[
W_n(x)=B_{0,n}(x)+B_{1,n}(x). \tag{22}
\]
Equations (14), (20)--(22), and \(h_\star\ge s_{n,\kappa}\) imply
\[
\begin{aligned}
\mathbb E_P\sup_{x\in\mathcal B}W_n(x)
&\le Cs_{n,\kappa}
+C\sqrt{\log n}\,s_{n,\kappa}^{-1}
\exp\{-cns_{n,\kappa}^\kappa\}\\
&\le Cs_{n,\kappa}.
\end{aligned}\tag{23}
\]
For the last inequality, uniformly over \(\kappa\in(2,\overline\kappa]\),
\[
ns_{n,\kappa}^\kappa
\ge n^{2/(\overline\kappa+2)}
(\log n)^{\overline\kappa/(\overline\kappa+2)},
\]
for all sufficiently large \(n\); hence the exponential in (23) dominates the displayed polynomial and logarithmic factors. Taking the supremum over \(P\in\mathcal P_{\kappa}^{\mathrm{perv}}\) proves the claimed pervasive expected-width rate, including the empty-admissible-set fallback through (21).

It remains to prove the general obstruction. Let a band have simultaneous coverage at least \(1-\gamma-\epsilon_n\) over a class containing \(P_0,P_1\), fix \(x\in\mathcal B\), and define
\[
A_j=\{\tau_{P_j}(x)\in\mathcal C_n(x)\},
\qquad j=0,1.
\]
Simultaneous coverage implies
\[
P_j^{\otimes n}(A_j)\ge1-\gamma-\epsilon_n,
\qquad j=0,1. \tag{24}
\]
By the definition of total variation,
\[
P_0^{\otimes n}(A_1)
\ge P_1^{\otimes n}(A_1)
-\operatorname{TV}(P_0^{\otimes n},P_1^{\otimes n}). \tag{25}
\]
The union bound under \(P_0^{\otimes n}\), followed by (24)--(25), gives
\[
P_0^{\otimes n}(A_0\cap A_1)
\ge1-2\gamma-2\epsilon_n
-\operatorname{TV}(P_0^{\otimes n},P_1^{\otimes n}). \tag{26}
\]
On \(A_0\cap A_1\), the interval \(\mathcal C_n(x)\) contains both target values. Therefore its diameter is at least their absolute difference, and \(\mathrm{def:band\text{-}half\text{-}width}\) gives
\[
W_n(x)\ge\frac12|\tau_{P_1}(x)-\tau_{P_0}(x)|. \tag{27}
\]
Multiplying (27) by the indicator of \(A_0\cap A_1\), taking \(P_0^{\otimes n}\)-expectations, and using the positive part of (26) proves
\[
\mathbb E_{P_0}W_n(x)\ge\frac{|\tau_{P_1}(x)-\tau_{P_0}(x)|}{2}
\left[1-2\gamma-2\epsilon_n-
\operatorname{TV}(P_0^{\otimes n},P_1^{\otimes n})\right]_+.
\]

This last inequality necessarily uses the total variation distance between the full observed experiments. Projection from the full data to the scores gives the opposite comparison,
\[
\operatorname{TV}(P_{0,X}^{\otimes n},P_{1,X}^{\otimes n})
\le\operatorname{TV}(P_0^{\otimes n},P_1^{\otimes n}),
\]
so score-law distance alone does not supply the upper bound on the full distance needed in (26). Likewise, (19) forces a global \(\log n\) calibration, while \(\mathrm{def:local\text{-}oracle\text{-}benchmark}\) may contain only bounded \(\log\{eK_P(x,h)/\gamma\}\) at an isolated thin site. The proof therefore establishes exactly the global honest band, its pervasive expected-width rate, and the full-experiment two-law obstruction asserted in the theorem, while leaving the expressly stated local-oracle ratio and exact score-only adaptation criterion unresolved.